\begin{frame}
Static Distance Estimation
\begin{itemize}
  \item Scenario:  Two soldiers each have a laser range finder.  The first soldier points his 
range finder at a potential target and determines it to be 205.78 yards away.  The second soldier 
determines with his range finder that the target is 201.22 yards away.
  \item How do these two soldiers combine their two measurements to 
produce an better estimate?
\end{itemize} 
\end{frame}

\begin{frame}
Static Distance Estimation
\begin{itemize}
  \item Soldier A says, "Let's just take the average of our two
measurements and call that in to the shooters."
  \item Soldier B disagrees and says, "No way!  My laser range finder
is much better than yours.  We should just use my measurrement as the 
the best estimate."
\end{itemize} 
\end{frame}


\begin{frame}
Static Distance Estimation
\begin{itemize}
  \item Actually, both soldiers are wrong!
  \item They should use a weighted average of both measurements as their final estimate.
  \item But what's the weighting scheme?
\end{itemize} 
\end{frame}

\begin{frame}
Static Distance Estimation
\begin{itemize}
  \item Intuitively, the weighting scheme should take into account the relative uncertainty of each measurement.
  \item Also, if the two weights are equal, the weighted average should reduce to the simple unweighted average.
  \item Let $r_a$ be the range measurment of Soldier A, $r_b$ that of Solider  B.  Then, a weighted average would look like this:
\[\hat{r}=w_a r_a+w_b r_b,\]
where $w_a$ is the weight associated with Solider A's measurement and $w_b$ is that associated with Solider B's measurement.
\end{itemize} 
\end{frame}

\begin{frame}
Static Distance Estimation
\begin{itemize}
  \item One way to set the weights using the uncertainty is to use the standard deviations associate with each measurement.
  \item Let $\sigma_a$ be the standard deviation associated with laser range finder of Soldier A, $\sigma_b$ that of Soldier B.
  \item Then we could set 
\[w_a=\dfrac{\sigma_b^2}{\sigma_a^2+\sigma_b^2}\;\;\;\text{and}\;\;\;
  w_b=\dfrac{\sigma_a^2}{\sigma_a^2+\sigma_b^2}.
\]
\end{itemize} 
\end{frame}

\begin{frame}
Static Distance Estimation
\begin{itemize}
  \item Why does $w_a=\dfrac{\sigma_b^2}{\sigma_a^2+\sigma_b^2}$ have $\sigma_b$ in its denominator instead of $\sigma_a$?
  \item Note also that these weights reduce to 1/2, as they should, when $\sigma_a=\sigma_b$.
\end{itemize} 
\end{frame}


\begin{frame}
Static Distance Estimation
\begin{itemize}
  \item If we have several sequential measurmeents, from a variety of range measurement sources, we would want to put this into a recursive form
like we did for the running average of the resistor.
  \item The recursive, weighted running average is beginning to look more and more like a Kalman filter.
\end{itemize} 
\end{frame}

\begin{frame}
Static Distance Estimation
\begin{itemize}
  \item To see this, we have to first realize that the information or certainty (=inverse of uncertainty) 
of the weighted average estimate is:
\[1/p=1/\sigma_a^2 + 1/\sigma_b^2.\]
 \item Note that $p$ is the scalar version of the estimation error covariance matrix $P$ and the $1/\sigma$ terms
are the measurement error covariances $R$.
\end{itemize} 
\end{frame}



\begin{frame}
Static Distance Estimation
\begin{itemize}
  \item This generalizes easily to $k$ measurements (let's number our soldiers now from 1 to $k$):
\[1/p(k)=1/\sigma_1 + \cdots + 1/\sigma_{k}=1/p(k-1)+1/\sigma_{k}\]
  \item The startling thing about this equation is that the infomration always increases, no matter how inaccurate the
measurement!  As long as the measurements are unbiased (i.e., zero mean) and Gaussian, it doesn't
 matter how inaccurate they are -- each new measurement will convey some amount of additional information
to improve our estimate.
\end{itemize} 
\end{frame}


\begin{frame}
Static Distance Estimation
\begin{itemize}
  \item Finally, our recursive estimate of range $\hat{r}(k)$ after the $k^\text{th}$ measurement is
\[\hat{r}(k) = \hat{r}(k-1)+ K(k)(r(k) - \hat{r}(k-1))\]
where
\[K(k) = p(k-1)\left[p(k-1)+\sigma_k^2\right]^{-1}\] and
\[p(k)=\left[1/p(k-1) + 1/\sigma_k^2\right]^{-1}\]
\end{itemize} 
\end{frame}

